% contributors: 
\tran{{\kofont 한국어} (ISO 639-3 kor)}{{\kofont 요네다 보조정리}}{{\kofont
	$\mathcal C$가 작은 범주이고 $F$가 이 범주에서 집합의 범주로 가는 함자라 하자. 그러면 $\mathcal C$의 모든 대상 $X$에 대하여 자연 변환의 집합 $\hom(-,X)\Rightarrow F$과 집합 $FX$ 사이에 자연스러운 전단사 함수가 존재하며, 이 동형 사상은 대응
	\[
		\Big(\xi : \hom(-,X)\Rightarrow F\Big) \mapsto \xi_X(1_X)
	\]
	에 의해 정의된다} ({\kofont 이 함수는 전단사이다}).
}
